\section{The model and how the engine is organised}

\subsection{The model in one page}

\paragraph{The game.} Individuals meet in pairs. In each meeting one is the \emph{Can} (the donor) and the other the
\emph{Need} (the receiver): each individual can help with probability $q$, a given ordered assignment of the two roles
therefore has probability $\kappa = q(1-q)$ (the \emph{role factor}, $1/4$ for $q=\tfrac12$). The donor chooses whether to
help. Helping costs the donor $c$ and gives the receiver $b$.

\paragraph{Scores.} Every individual has an integer score $j\in\{0,1,2,\dots\}$ (a reputation; there is no cap).
After a meeting the \emph{scoring norm} changes the donor's score, and the receiver's, by $-1$, $0$ or $+1$, depending on
whether help was given and on the \emph{classes} of the two scores. A class is a set of scores: for binary norms
$s=0$ and $s\ge 1$, for ternary norms $s=0$, $0<s\le\tau$ and $s>\tau$. A score of $0$ cannot go down.

\paragraph{Strategies.} A strategy says with which probability a donor helps. It may depend on the receiver's class
and on the donor's own score. A population holds $n$ strategies with frequencies $\rho_i$ (summing to 1). Each strategy
$i$ has its own score distribution $w^{(i)}_j$ = the fraction of its individuals with score $j$.

\paragraph{Score chains.} Seen from one individual of strategy $i$, the score is a birth--death chain on $\{0,1,\dots\}$:
it moves up from $j-1$ to $j$ at rate $g^{(i)}_j$ and down from $j$ to $j-1$ at rate $f^{(i)}_j$. The rates come from
the norm, the strategies, who is in the population and, through noise $\varepsilon$, from random score changes. They
depend on the score distributions of \emph{all} strategies, so the $n$ chains are run together until the distributions
stop changing (the \emph{stable} score distributions).

\paragraph{Payoffs, replicator dynamics, ESS.} With the stable distributions, the probability that a Can of
strategy $i$ helps a Need of strategy $k$ is $H_{ik}$, and the expected payoff of strategy $i$ is
\[
 \pi_i \;=\; \kappa\Bigl(b\sum_k \rho_k H_{ki} \;-\; c\sum_k \rho_k H_{ik}\Bigr).
\]
The frequencies then evolve by the replicator equation $\dot\rho_i=\rho_i(\pi_i-\bar\pi)$, and a strategy is an
\emph{evolutionarily stable strategy} (ESS) if a rare invader cannot do better than it.

\subsection{Notation and the names in the code}

\begin{center}
\begin{tabularx}{\linewidth}{@{}llX@{}}\toprule
symbol & in the code & meaning\\\midrule
$j$ & \code{j}, a column of \code{w} & a score, $0,1,2,\dots$\\
$i,k$ & a row of \code{w}, \code{rho} & a strategy\\
$\rho_i$ & \code{rho[i]} & frequency of strategy $i$\\
$w^{(i)}_j$ & \code{w[i, j]} & fraction of strategy $i$ with score $j$ (array of shape $n\times L$)\\
$\kappa$ & \code{role\_factor} & $q(1-q)$, default $0.25$\\
$G^{(i)}_j,\;F^{(i)}_j$ & \code{G[i, j]}, \code{F[i, j]} & rate $j-1\to j$ and $j\to j-1$ from interactions alone\\
$g^{(i)}_j,\;f^{(i)}_j$ & \code{g[i, j]}, \code{f[i, j]} & the same with noise\\
$\varepsilon,\varepsilon_\uparrow,\varepsilon_\downarrow$ & \code{eps}, \code{eps\_up}, \code{eps\_down} & noise\\
$h_i(c,d)$ & \code{donor\_help(n)[c, d]} & probability that a donor of strategy $i$ with score $d$ helps a receiver of class $c$\\
$H_{ik}$ & \code{H[i, k]} & probability that a Can of $i$ helps a Need of $k$\\
$A_{ik}$ & \code{A[i, k]} & payoff matrix, $\pi=A\rho$\\
$\pi_i$ & \code{payoffs[i]} & expected payoff of strategy $i$\\
$b,\;c$ & \code{benefit}, \code{cost} & benefit to the receiver, cost to the donor\\
\bottomrule
\end{tabularx}
\end{center}

Arrays are indexed from $0$. The rate arrays have one unused column: \code{g[:, 0]} and \code{f[:, 0]} are $0$, and
\code{g[i, j]} (\code{G[i, j]}) is the rate \emph{into} score $j$ (so $j\ge1$), while \code{f[i, j]} (\code{F[i, j]}) is the
rate \emph{out of} score $j$ (also $j\ge1$). For $L$ scores, \code{g} has $L+1$ columns and \code{f} has $L$.

\subsection{How the pieces fit together}

\begin{center}
\begin{tabularx}{\linewidth}{@{}lXl@{}}\toprule
chapter & pieces (module) & builds on\\\midrule
\ref{ch:elements} Elements & score partitions (\code{partitions}), strategies (\code{strategies}), scoring norms and noise (\code{norms}, \code{known\_norms}, \code{noise}) & ---\\
\ref{ch:scores} Score dynamics & population (\code{population}), transition rates (\code{rates}), the chain (\code{chain}), stable scores (\code{solver}) & elements\\
\ref{ch:dynamics} Dynamics & payoffs (\code{payoffs}), replicator dynamics (\code{replicator}), the strategy simplex and rest points (\code{simplex}), ESS (\code{ess}) & score dynamics\\
\bottomrule
\end{tabularx}
\end{center}

Everything is imported from the package: \code{from engine import PhiDisc, ScoringNorm, \dots}. Every example in this
document is run when the document is built, and what is printed below it is what the engine printed.

\paragraph{How each entry is laid out.} \emph{What it does}, \emph{how it does it}, its \emph{parameters} (name, type
with the default, meaning), and an \emph{example}.
